% Limitations — required by ACL/EMNLP

\paragraph{Architecture and scale.}
Our primary results are on DeepSeek-R1-Distill-Qwen-7B~\citep{deepseek2025r1} with GQA ($G{=}7$).
Cross-architecture validation on Qwen2.5-Math-7B covers SFT recovery only ($N{=}200$); both models share the Qwen2.5 base architecture, so our cross-architecture evidence does not extend to different architecture families (e.g., LLaMA, Mistral).
Whether the spectrum generalizes across architecture families and scales (1B--70B) remains open.

\paragraph{Diagnostic limitations.}
Reward-accuracy decoupling at 20--25\% sparsity means that recovery capacity---operationalized as peak reward---does not predict downstream accuracy improvements within our training horizon.
Recovery capacity captures the model's ability to sustain RL training without collapse, but practitioners seeking accuracy gains may require longer training, different reward formulations, or lower sparsity levels.
The per-layer capacity floor (40\% retention) was chosen heuristically and not ablated; different floor values may shift the recovery capacity spectrum.

\paragraph{Replication and scope.}
CKA-only and Standard are validated over 3 GRPO seeds; \rasp{} and Random each use a single seed ($N{=}1$).
We evaluate on four benchmarks (GSM8K, MATH-500, HumanEval, MBPP) covering math reasoning and code generation; logical, commonsense, and multi-modal reasoning are not addressed.
